

\documentclass[11pt]{article}
\usepackage{amsmath, amssymb, amsthm}
\usepackage{geometry}
\geometry{margin=1in}
\title{Foundations of Machine Learning -- Lecture 4 Notes}
\author{}
\date{}

\begin{document}
\maketitle

\section*{Classification}
\textbf{Method:} Pose the problem as linear regression

\[
	f(x) =
	\begin{cases}
		1 & w^T \hat{x} > 0.5    \\
		0 & w^T \hat{x} \leq 0.5
	\end{cases}
\]

We have a problem though that this function is not differentiable!
So the solution is to use a sigmoid function:

\[
	\sigma(t) = \frac{1}{1 + e^{-t}}
\]

\[
	f(x) = \sigma(w^T \hat{x}) = \frac{1}{1 + e^{-w^T \hat{x}}}
\]


\subsection*{Log-Likelihood}

\begin{itemize}


	\item Using MSE leads to a non convex function
	\item The labels that we are predicting are binary: $y \in \{0,1\}$.
	\item Meaning the prediction $f(x)$ can be thought of as the probability output
	\item Given the probability $f(x)$ what is the likelihood of passing $(y=1)$? What is the likelihood of failing $(y=0)$?
\end{itemize}

\[
	p(y|x,w) = f(x)^y (1 - f(x))^{(1-y)}
\]


\[
	p(y=1|f(x)) = f(x), \quad p(y=0|f(x)) = 1 - f(x)
\]


Instead of writing two separate cases for $y=0$ and $y=1$ we can combine them into one compact expression. Because the exponent $y$ "selects" the correct term:
If $y=1$ the first term survives.
If $y=0$ the second term survives.

\pagebreak

For observations $\{(x^{(i)}, y^{(i)})\}$ assuming i.i.d. we have:

\[
	likelihood = \prod_{i=1}^N f(x^{(i)})^{y^{(i)}} \big(1 - f(x^{(i)})\big)^{(1 - y^{(i)})}
\]

\[
	log\text{--}likelihood = - \sum_{i=1}^N \Big[ y^{(i)} \log \big(f(x^{(i)})\big)
		+ \big(1 - y^{(i)}\big) \log \big(1 - f(x^{(i)})\big) \Big]
\]

Therefore,

\[
	\mathcal{L}(w) = - \sum_{i=1}^{N} \left[ y^{(i)} \log\!\big(f(x^{(i)})\big) + \big(1 - y^{(i)}\big)\log\!\big(1 - f(x^{(i)})\big) \right]
\]

\[
	\nabla_w \mathcal{L}(w) = \sum_{i=1}^N x^{(i)} \Big(f(x^{(i)}) - y^{(i)}\Big)
\]

\[
	\nabla_w^2 \mathcal{L}(w) = \sum_{i=1}^N x^{(i)} {x^{(i)}}^T f(x^{(i)}) \Big(1 - f(x^{(i)})\Big)
\]

\subsection*{Multi-Class Classification}

In multi-class settings with $K > 2$, we reduce the problem to multiple binary classification tasks. Two common strategies are \textbf{One-vs-All} and \textbf{One-vs-One}.

\subsubsection*{One-vs-All}
Train $K$ binary classifiers.
Classifier $f^{(i)}(x)$ learns to distinguish class $i$ from all other classes.

\[
	\hat{y} = \arg\max_{i \in \{1,\dots,K\}} f^{(i)}(x)
\]

\begin{itemize}
	\item Simple to implement; only $K$ classifiers.
	\item Each model solves an imbalanced problem (one class vs.\ all others).
	\item Boundaries can be harder to learn because each classifier must model $K-1$ negative classes.
\end{itemize}

\subsubsection*{One-vs-One}
Train a classifier for every pair of classes, giving
\[
	\frac{K(K-1)}{2} \text{ classifiers}.
\]

Each classifier votes for one of its two classes.
The predicted label is the majority vote.

\begin{itemize}
	\item Each classifier only sees 2 classes, giving simpler decision boundaries and cheaper evaluation per model.
	\item Requires many classifiers, scaling as $O(K^2)$.
	\item Voting ties can occur and require tie-breaking.
\end{itemize}

\end{document}

